% Lösungen Lineare Funktionen · Anwendung · Prüfungs-Fokus MSA (zusammenbau v0.6)
\einheitenkopf{Einheit 5 · Anwendungen}

\erg{1}{a) Anfangswert 45 € ($n = 45$), Änderung 38 € je Stunde ($m = 38$) \quad b) $y = 3x + 8$ \quad c) $m = 6$, $n = 10$: $y = 6x + 10$}
\erg{2}{a) Angebot 1: $4 \cdot 32 + 120 \cdot 0{,}30 = 164$ €; Angebot 2: $95 + 420 \cdot 0{,}12 = 145{,}40$ €; Angebot 2 ist günstiger; $y = 0{,}12x + 95$ \quad b) Angebot 1: $3 \cdot 29 + 130 \cdot 0{,}40 = 139$ €; Angebot 2: $110 + 380 \cdot 0{,}10 = 148$ €; Angebot 1 ist günstiger; $y = 0{,}10x + 110$ \quad c) 240 km: Nord $150 + 1{,}20 \cdot 240 = 438$ €, Süd $1{,}80 \cdot 240 = 432$ €, Süd günstiger; 320 km: Nord $150 + 1{,}20 \cdot 320 = 534$ €, Süd $1{,}80 \cdot 320 = 576$ €, jetzt Nord günstiger (gleich teuer bei 250 km) \quad d) 260 km: Adler $90 + 0{,}80 \cdot 260 = 298$ €, Biber $1{,}10 \cdot 260 = 286$ €, Biber günstiger; 320 km: Adler $90 + 0{,}80 \cdot 320 = 346$ €, Biber $1{,}10 \cdot 320 = 352$ €, jetzt Adler günstiger (gleich teuer bei 300 km) \quad e) erster Sachverhalt: $y = 0{,}25x + 1$ (Cent in Euro umrechnen), zweiter: $y = x + 25$ \quad f) erster Sachverhalt: $y = 0{,}60x + 2$ (Cent in Euro umrechnen), zweiter: $y = 2x + 60$ \quad g) $m = 45$, $n = 640$: $y = 45x + 640$; aus $45x + 640 = 1\,500$ folgt $x \approx 19{,}1$, also nach 20 Monaten \quad h) $m = 70$, $n = 1\,150$: $y = 70x + 1\,150$; aus $70x + 1\,150 = 2\,000$ folgt $x \approx 12{,}1$, also nach 13 Monaten \quad i) $m = 1{,}2$, $n = 180$: $y = 1{,}2x + 180$ \quad j) $m = 2{,}4$, $n = 250$: $y = 2{,}4x + 250$ \quad k) $830 + 12 \cdot 35 = 1\,250$ € \quad l) $2\,340 + 6 \cdot 85 = 2\,850$ Bäume \quad m) zu Beginn $300$ l (30 l mehr als nach 5 min); nach 40 min: $300 - 8 \cdot 30 = 60$ l \quad n) zu Beginn $20$ \% (6 Prozentpunkte weniger als nach 10 min); nach 90 min: $20 + 9 \cdot 6 = 74$ \% \quad o) Blitz: II, Flott: I; II beginnt bei 40 € auf der y-Achse – das ist der Grundpreis \quad p) Hoch: I, Weit: II; I beginnt bei 30 € auf der y-Achse – das ist die Aufnahmegebühr \quad q) erstes Angebot: A (bis zu den Freikilometern gleichbleibend, dann steigend); zweites Angebot: C (beginnt bei der Einmalzahlung, steigt gleichmäßig) \quad r) erstes Angebot: D (bis zu den Freikilometern gleichbleibend, dann steigend); zweites Angebot: B (beginnt bei der Einmalzahlung, steigt gleichmäßig)}
